\chapter{Several GPUs}\label{ch:mgpu}

The port runs one domain on one GPU. Two reasons call for more: the full case does not fit a 40~GB GPU, which is what
the project's workstation has (two A100 40~GB), and larger fires need larger domains. WRF already runs on many CPUs
by splitting each domain into patches, one per MPI rank, and exchanging halos (\cref{ch:cpupar}). The multi-GPU port
keeps that design and moves the exchanges onto the GPUs. It is Stage~8 of the plan (phases S8-01 to S8-05).

\section{Halo exchanges in WRF}

Each patch keeps a halo: copies of the neighbouring patches' points that its stencils read. Before a stencil needs
them, the halo is refreshed. The Registry generates one include file per halo (\code{HALO_EM_A.inc}, \dots), naming the
fields and the halo width, and RSL\_LITE does the work (\code{external/RSL_LITE}):
\begin{enumerate}
  \item \code{RSL_LITE_INIT_EXCH} sizes the buffers for the neighbours;
  \item \code{RSL_LITE_PACK} copies each field's edge rows into a send buffer, through the loops of \code{f_pack.F90};
  \item \code{RSL_LITE_EXCH_Y} and \code{RSL_LITE_EXCH_X} exchange the buffers with the neighbours in $y$, then in
    $x$, with non-blocking \code{MPI_Isend} and \code{MPI_Irecv};
  \item the same pack routine, in unpack mode, copies the received buffers into the halos.
\end{enumerate}
The $y$ exchange comes first and the $x$ exchange second, so that the corners arrive through two hops.

\section{Halos on the device}

With the state on the GPU, the pack and unpack loops must run on the GPU too; otherwise every exchange would copy the
whole patch to the host. The design (phase S8-01, an ADR to write) has three parts:
\begin{itemize}
  \item \textbf{Device pack and unpack.} The loops of \code{f_pack.F90} become kernels writing a device buffer.
  \item \textbf{GPU-aware MPI.} An MPI library built with CUDA support accepts device addresses in \code{MPI_Isend}
    and \code{MPI_Irecv}. The OpenACC device address of the buffer is passed with \code{host_data use_device}. Between
    GPUs of one node, the transfer can then go directly over NVLink or PCIe, without passing through host memory.
  \item \textbf{Generated halo includes on the device.} The halo includes generated by the Registry call the device
    versions when the domain lives on the device.
\end{itemize}
Packing, sending and unpacking copy values; they change no arithmetic (class R0). What can change results is whether
the computation on each patch is the same as on the whole domain.

\section{Bitwise across decompositions}

A one-GPU run and a two-GPU run of the same case should give the same bits. WRF is designed for this: every point is
computed from its own stencil, whichever patch it lies in, and halos supply exactly the values the whole domain would
have. The exceptions are the places where results depend on the decomposition:
\begin{itemize}
  \item loops whose bounds come from the patch where they should come from the domain;
  \item reductions across patches: sums over the domain added in a different order. Maxima and integer counts are
    exact; floating-point sums are not.
\end{itemize}
CPU-REF itself must be decomposition-independent, which test T-DEC checks before the reference runs: 1 rank against
many on the first steps, and different rank counts on 30 minutes (\code{plan.md} P0.10). On the GPUs, T-DEC-GPU
compares one GPU with two on the dev case (phase S8-03). A failure is bisected with the tracer to the first routine
whose result depends on the decomposition, and that routine is fixed in both builds so that its result equals the
one-rank result.

\section{The full case on two A100 40 GB}

The full case needs about 57~GB (\cref{ch:perf}). Split over two GPUs, the state of each domain halves, but some
items do not:
\begin{itemize}
  \item halved: the state of d01 and d02 including the fire grid (31.4~GB), the scratch pool and the work arrays;
  \item not halved: the RRTMG batch (4.5~GB, unless the batch is made smaller), the local-memory reservation of
    column physics (proportional to the GPU's resident threads, 6.6~GB on an A100), the CUDA context;
  \item added: halos and exchange buffers.
\end{itemize}
The estimate per GPU and the measured peak are \tbr\ (phase S8-04). Once the case fits, the workstation can run the
acceptance criteria of G5 itself.

\section{Scaling}

A patch's work grows with its area and its halo traffic with its perimeter. For a domain of $n \times n$ columns split
into $p$ patches in a square layout, each patch has about $n^2/p$ columns and exchanges about $4\,w\,n/\sqrt{p}$ halo
columns of width $w$ per exchange. The ratio of exchange to work grows like $\sqrt{p}/n$. For d02 ($n = 810$) on two
GPUs it is small; the question is whether the many exchanges per step (several per RK stage and per acoustic
sub-step) add latency. Two remedies keep the bits:
\begin{itemize}
  \item overlap: compute the interior of the patch while the halo is in flight, then the edge points. Every point is
    computed with the same statements; only the order in which different points are computed changes, and points do
    not depend on each other within a kernel (class R0);
  \item fewer, wider exchanges: exchange a wider halo less often, recomputing the extra points with the same
    operations (class R1).
\end{itemize}
The measured scaling from one to two GPUs, on the A100 workstation and on an H100 node, is \tbr\ (phase S8-05).
